\subsection{Two native clocks and a controlled common-clock limit}
The actual eight-level native energy admits an exact Jordan--Wigner representation as three normal fermion occupations. Two occupations therefore supply independently prepared clock bits within the same source cell. Their joint Ramsey instrument retains all settings, both clock outputs and the spectator complement. They share one fixed quantum geometry and scalar-wave source of dimension $80000$.

Complete records resolve a finite response difference. At event 16, the difference between the two background-induced phases after normalization by their respective bare gaps lies in
\[
[3.06881,\ 3.21525]\times10^{-7},
\]
corresponding to $31.37$--$32.87$ parts per million relative to the second clock response. An exact dyadic partial-transpose witness in the vacuum-field controls also certifies correlations that entangle the clocks through their common geometry. Mean wave-arrival delays measured by both clocks are positive and use the same receiver law. These finite source effects are distinguished from a measured equivalence-principle violation.

There is nevertheless a controlled common-clock limit. For the declared probe family $H_\lambda=H_{\rm env}+\lambda B(a n_a+c n_c)$, with $B=\exp(Q_{\rm right})$, both normalized clock phases converge to the same integral of the Heisenberg-evolved $B$. For $0<\lambda\leq10^{-3}$ and source parameter at most $8$, their disagreement is bounded by $1500\lambda/(1-600\lambda^2)$, uniformly over arbitrary environment states and passive references. This follows from the exact finite Duhamel remainder and the complete joint instrument. No physical time or gravitational constant calibration is assumed.
